\documentclass[11pt]{article}
\usepackage[margin=1in]{geometry}
\usepackage[T1]{fontenc}
\usepackage[utf8]{inputenc}
\usepackage{lmodern}
\usepackage{xcolor}
\usepackage{booktabs}
\usepackage{longtable}
\usepackage{array}
\usepackage{amsmath,amssymb}
\usepackage{enumitem}
\usepackage{hyperref}
\hypersetup{colorlinks=true, linkcolor=blue, urlcolor=blue}
\definecolor{accent}{RGB}{25,56,99}

\setlength{\parskip}{0.45em}
\setlength{\parindent}{0pt}

\begin{document}

\begin{center}
{\LARGE\bfseries\color{accent} Comparison of Power-Law and Logistic Models for SIP Saturation Relationships}\\[0.4em]
{\large Statistical comparison of saturation--electrical property models using the SIP analysis pipeline}
\end{center}

\section{Introduction}

Spectral Induced Polarization (SIP) measurements provide valuable information about subsurface electrical properties and their dependence on water saturation. Laboratory SIP column experiments often measure how electrical resistivity and conductivity evolve as the saturation state of the porous medium changes.

A key objective in analyzing SIP data is identifying the functional relationship between saturation ($S$) and electrical properties such as resistivity ($\rho$) and conductivity ($\sigma$). This report compares two commonly used model families:

\begin{itemize}[leftmargin=1.5em]

\item power-law scaling models

\item logistic (sigmoidal) models

\end{itemize}

The analysis is based on datasets processed using the SIP pipeline, which extracts and merges experimental measurements from multiple runs. Model comparison is performed using several statistical metrics including coefficient of determination ($R^2$), sum of squared errors (SSE), Gaussian log-likelihood, and Akaike Information Criterion (AIC).

\section{Dataset structure}

The comparison framework operates on datasets generated by the SIP processing workflow. Each dataset contains:

\begin{itemize}[leftmargin=1.5em]

\item water saturation ($S$)

\item electrical resistivity ($\rho$)

\item electrical conductivity ($\sigma$)

\item transformed variables (logarithmic or normalized forms)

\end{itemize}

Data may originate from multiple experimental locations, including:

\begin{itemize}[leftmargin=1.5em]

\item Rosemead

\item Van Nuys (VN)

\end{itemize}

The analysis can be performed at three hierarchical levels:

\begin{longtable}{p{0.22\textwidth} p{0.60\textwidth}}
\toprule
\textbf{Level} & \textbf{Description}\\
\midrule
\endhead
Sample & Each individual experimental run is analyzed separately.\\
Site & All runs from the same experimental site are pooled together.\\
All & The entire dataset is combined into a single regression dataset.\\
\bottomrule
\end{longtable}

This structure allows investigation of both local experimental variability and broader site-scale trends.

\section{Model formulations}

\subsection{Power-law model}

The power-law model is commonly used in petrophysical relationships and has the form:

\[
y = a S^{b}
\]

where:

\begin{itemize}

\item $S$ = water saturation

\item $a$ = scaling coefficient

\item $b$ = saturation exponent

\end{itemize}

Taking the logarithm yields:

\[
\log y = b\,\log S + \log a
\]

This transformation converts the relationship into a linear regression problem in log–log space.

The pipeline supports both base-10 and natural logarithm transformations:

\[
\log_{10}(S) \quad \text{or} \quad \ln(S)
\]

depending on the selected configuration.

\subsection{Logistic model}

To capture nonlinear transitions, a four-parameter logistic model is used:

\[
y = c + \frac{L}{1 + \exp[-k(S-x_0)]}
\]

where:

\begin{itemize}

\item $L$ = amplitude

\item $k$ = steepness of the transition

\item $x_0$ = transition saturation

\item $c$ = baseline offset

\end{itemize}

The logistic function can represent threshold-driven responses where electrical properties change rapidly within a limited saturation range.

In the current implementation, the amplitude parameter may be positive or negative, allowing the model to represent both increasing and decreasing responses.

\section{Statistical evaluation}

To compare candidate models, several statistical metrics are computed.

\subsection{Sum of squared errors}

\[
SSE = \sum (y_i - \hat{y}_i)^2
\]

Lower SSE values indicate a closer match between predicted and observed values.

\subsection{Coefficient of determination}

\[
R^2 = 1 - \frac{SSE}{SST}
\]

where

\[
SST = \sum (y_i - \bar{y})^2
\]

Higher values of $R^2$ indicate stronger explanatory power.

\subsection{Residual standard deviation}

The residual standard deviation provides a measure of the spread of model residuals:

\[
\sigma = \sqrt{\frac{SSE}{n}}
\]

\subsection{Log-likelihood}

Assuming Gaussian residuals, the log-likelihood is:

\[
\ln L =
-\frac{n}{2}
\left[
\ln(2\pi\sigma^2) + 1
\right]
\]

Larger values correspond to better fits.

\subsection{Akaike Information Criterion}

Model selection is primarily based on AIC:

\[
AIC = n\ln\left(\frac{SSE}{n}\right) + 2k
\]

where $k$ is the number of model parameters.

\begin{itemize}

\item Power-law model: $k=2$

\item Logistic model: $k=4$

\end{itemize}

The model with the smaller AIC is considered statistically preferable.

\section{Normalization strategy}

Optional normalization can be applied to both dependent and independent variables before fitting.

Normalization is always applied within the relevant analysis group:

\begin{itemize}

\item sample normalization uses only sample data

\item site normalization uses all runs from that site

\item global normalization uses the full dataset

\end{itemize}

This prevents inconsistencies that would occur if sample-level normalization were reused for site-level or global analyses.

Two normalization methods are available:

\subsection{Z-score normalization}

\[
x' = \frac{x-\mu}{\sigma_x}
\]

\subsection{Min–max normalization}

\[
x' = \frac{x-x_{\min}}{x_{\max}-x_{\min}}
\]

Normalization affects only the scale of variables and does not change the underlying relationships.

\section{Interpretation of model behavior}

The two model families represent different physical mechanisms.

\subsection{Power-law scaling}

Power-law relationships typically indicate continuous scaling behavior across the saturation range. Resistivity often follows this pattern because electrical conduction primarily occurs through the pore fluid phase.

\subsection{Logistic transition behavior}

Logistic relationships indicate the presence of threshold behavior. For conductivity, this may correspond to the formation of connected water pathways within the pore network.

Below a critical saturation level, conductive pathways remain disconnected and conductivity remains low. As saturation increases and connectivity improves, conductivity can increase rapidly before eventually stabilizing.

Such behavior is consistent with percolation theory and pore-network connectivity models.

\section{Typical trends}

Across many SIP experiments, the following general patterns are often observed:

\begin{longtable}{p{0.32\textwidth} p{0.55\textwidth}}
\toprule
\textbf{Relationship} & \textbf{Typical behavior}\\
\midrule
\endhead
Resistivity vs saturation & Often well described by power-law scaling.\\
Conductivity vs saturation & Frequently exhibits sigmoidal or threshold-like behavior.\\
Log-transformed resistivity & Often appears linear in log–log space.\\
Log-transformed conductivity & May remain nonlinear depending on site conditions.\\
\bottomrule
\end{longtable}

These trends can vary depending on soil type, pore structure, and experimental conditions.

\section{Output files}

The comparison workflow generates several output files summarizing model performance.

\begin{longtable}{p{0.40\textwidth} p{0.18\textwidth} p{0.28\textwidth}}
\toprule
\textbf{File} & \textbf{Scope} & \textbf{Purpose}\\
\midrule
\endhead
\texttt{compare\_power\_vs\_logistic.csv} & Global & Table containing regression statistics for each model.\\
\texttt{compare\_power\_vs\_logistic.xlsx} & Global & Excel summary with grouped results.\\
\texttt{stacked\_compare\_input.csv} & Campaign & Combined dataset used for fitting.\\
PNG plots & Per case & Visual comparison of model fits.\\
\bottomrule
\end{longtable}

Each output record includes:

\begin{itemize}

\item model parameters

\item $R^2$, SSE, $\sigma$, and $\ln L$

\item AIC values

\item model winner

\item analysis level (sample, site, or global)

\end{itemize}

\section{Conclusion}

This report framework provides a systematic comparison between power-law and logistic models for describing saturation–electrical property relationships in SIP experiments.

Key observations include:

\begin{itemize}

\item resistivity commonly follows power-law scaling

\item conductivity often exhibits nonlinear or threshold behavior

\item logistic models are useful for capturing rapid transitions associated with pore connectivity

\end{itemize}

By combining statistical model comparison with physically interpretable functions, the SIP pipeline enables more robust interpretation of laboratory SIP experiments.

\end{document}
